\documentclass[../../main]{subfiles}

\begin{document}

\section{Fundamentos}
\label{sec:matematica:teoria-conjuntos:fundamentos}

\begin{defn}[Conjunto]
	\label{def:matematica:teoria-conjuntos:fundamentos:conjunto}

	Um conjunto é um conceito primitivo que representa uma coleção bem
	definida de objetos, denominados elementos do conjunto.
\end{defn}

\begin{defn}[Elemento]
	\label{def:matematica:teoria-conjuntos:fundamentos:elemento}

	Um elemento é um objeto pertencente a um conjunto.
\end{defn}

\begin{defn}[Pertinência]
	\label{def:matematica:teoria-conjuntos:fundamentos:pertinencia}

	A pertinência é a relação entre um elemento e um conjunto.

	Dados um elemento \(x\) e um conjunto \(A\),

	\[
		x \in A
	\]

	significa que \(x\) pertence a \(A\).

	Analogamente,

	\[
		x \notin A
	\]

	significa que \(x\) não pertence a \(A\).
\end{defn}

\begin{defn}[Inclusão]
	\label{def:matematica:teoria-conjuntos:fundamentos:inclusao}

	Sejam \(A\) e \(B\) conjuntos.

	Dizemos que \(A\) está contido em \(B\), ou que \(A\) é um subconjunto
	de \(B\), quando todo elemento de \(A\) também pertence a \(B\).

	Formalmente,

	\[
		A \subseteq B
		\iff
		(\forall x)
		(x \in A \Rightarrow x \in B).
	\]
\end{defn}

\begin{defn}[Subconjunto Próprio]
	\label{def:matematica:teoria-conjuntos:fundamentos:subconjunto-proprio}

	Sejam \(A\) e \(B\) conjuntos.

	Dizemos que \(A\) é um subconjunto próprio de \(B\) quando

	\[
		A \subseteq B
	\]

	e

	\[
		A \neq B.
	\]

	Neste caso escrevemos

	\[
		A \subset B.
	\]
\end{defn}

\begin{defn}[Igualdade de Conjuntos]
	\label{def:matematica:teoria-conjuntos:fundamentos:igualdade-conjuntos}

	Dois conjuntos são iguais quando possuem exatamente os mesmos
	elementos.

	Formalmente,

	\[
		A = B
		\iff
		(\forall x)
		(x \in A \Leftrightarrow x \in B).
	\]
\end{defn}

\begin{prop}[Princípio da Extensionalidade]
	\label{prop:matematica:teoria-conjuntos:fundamentos:extensionalidade}

	Dois conjuntos são iguais se, e somente se, possuem exatamente os
	mesmos elementos.

	Em símbolos,

	\[
		(\forall x)
		(x \in A \Leftrightarrow x \in B)
		\Rightarrow
		A = B.
	\]
\end{prop}

\begin{proof}

	A afirmação decorre diretamente da definição de igualdade de
	conjuntos.

\end{proof}

\begin{defn}[Conjunto Vazio]
	\label{def:matematica:teoria-conjuntos:fundamentos:conjunto-vazio}

	O conjunto vazio é o conjunto que não possui elementos.

	É representado por

	\[
		\varnothing.
	\]

	Formalmente,

	\[
		(\forall x)
		(x \notin \varnothing).
	\]
\end{defn}

\begin{prop}[Unicidade do Conjunto Vazio]
	\label{prop:matematica:teoria-conjuntos:fundamentos:unicidade-conjunto-vazio}

	Existe um único conjunto vazio.
\end{prop}

\begin{proof}

	Sejam \(A\) e \(B\) conjuntos vazios.

	Como nenhum elemento pertence a \(A\) e nenhum elemento pertence a
	\(B\),

	\[
		(\forall x)
		(x \in A \Leftrightarrow x \in B).
	\]

	Pelo princípio da extensionalidade,

	\[
		A = B.
	\]

	Portanto o conjunto vazio é único.

\end{proof}

\begin{prop}
	\label{prop:matematica:teoria-conjuntos:fundamentos:vazio-subconjunto}

	O conjunto vazio é subconjunto de qualquer conjunto.
\end{prop}

\begin{proof}

	Seja \(A\) um conjunto.

	Deve-se mostrar que

	\[
		(\forall x)
		(x \in \varnothing \Rightarrow x \in A).
	\]

	Como não existem elementos em \(\varnothing\), a implicação é
	verdadeira para todo \(x\).

	Logo,

	\[
		\varnothing \subseteq A.
	\]

\end{proof}

\begin{defn}[Descrição por Extensão]
	\label{def:matematica:teoria-conjuntos:fundamentos:extensao}

	Um conjunto pode ser descrito pela enumeração explícita de seus
	elementos.

	Exemplo:

	\[
		A = \{1,2,3\}.
	\]
\end{defn}

\begin{defn}[Descrição por Compreensão]
	\label{def:matematica:teoria-conjuntos:fundamentos:compreensao}

	Um conjunto pode ser descrito por uma propriedade satisfeita por seus
	elementos.

	Exemplo:

	\[
		A =
		\{
		x \in \mathbb{N}
		\mid
		x < 4
		\}.
	\]
\end{defn}

\end{document}
